\chapter{Label Construction Rules}
\label{app:labels}

\section{OpenLane-V2}

The label builder (\code{dtsu/labels.py}) processes each frame as follows.

\begin{enumerate}
  \item \textbf{Traffic elements.} Read the front-image traffic elements and their attribute: 1 red,
        2 green, 3 yellow for circular lights; 4--12 for arrow types (left, right, straight, U-turn and
        their combinations).
  \item \textbf{Ego approach lanes.} Select lane centerlines that start ahead of the ego vehicle and are
        aligned with its heading (within $60^\circ$).
  \item \textbf{Maneuver per lane.} From each approach lane, follow \code{topology_lclc} to the first
        junction lane being entered; the entering lane may start up to 3\,m behind the ego vehicle
        (\code{ENTERING_X = -3 m}). Classify the heading change along that lane: $|\Delta\psi|\le30^\circ$
        straight, $|\Delta\psi|>135^\circ$ U-turn, otherwise left or right by the sign of $\Delta\psi$.
  \item \textbf{Relevant lights.} Through \code{topology_lcte}, collect the lights linked to the selected
        lanes of the nearest junction (within 10\,m). If the ego vehicle has passed the stop line, use
        approach lanes ending within 30\,m behind it.
  \item \textbf{Assignment.} For each relevant light, record the set of maneuvers it controls. Then:
        straight $\leftarrow$ straight-specific or generic lights; left $\leftarrow$ left-specific lights,
        else straight; U-turn $\leftarrow$ U-turn-specific lights, else left; right $\leftarrow$ lights
        that control only the right turn, else \emph{none}. Among candidates, the most specific light
        (fewest controlled maneuvers) wins.
  \item \textbf{Conflicts and filter.} Equally specific candidates with different colors mark a conflict
        and drop the frame. Frames without any relevant colored light are dropped.
  \item \textbf{Traffic side.} US (right-hand) rules are applied to all frames, including Singapore. A
        mirrored variant for left-hand traffic is kept for analysis only.
\end{enumerate}
Each record stores the image path, the four colors and passabilities, the traffic side, the
right-turn convention (\code{near_turn}) and all light boxes with their controlled maneuvers and a
relevance flag (used by the structured baseline).

\section{DTLD}

\begin{enumerate}
  \item Consider front-facing vehicle lamps (circles and arrows); pedestrian and other lamps are
        ignored.
  \item The ego intersection's lamps are the relevant ones plus non-relevant lamps at least 0.6 times
        the median height of the relevant ones.
  \item A route with its own arrow (left, right, straight, straight-left) is served by that arrow; a
        route without an arrow is served by the circles.
  \item States: red and red--yellow $\to$ red, yellow $\to$ yellow, green $\to$ green; off or unknown
        lamps give no color, and a direction whose lamps all lack one drops the frame.
  \item Directions: straight, left and right take their lamps' colors; left without lamps inherits
        straight; U-turn inherits left (DTLD has no U-turn signals). Right uses the German convention
        (\code{near_turn = "signal"}).
\end{enumerate}

\section{LISA}

\begin{enumerate}
  \item Lights are grouped by kind (circle, left arrow, straight arrow).
  \item A light is at the ego intersection if its height is at least 0.6 times that of the tallest light
        of the same kind in the frame.
  \item Directions follow the OpenLane-V2 rules under the US convention; right is always \emph{none}.
  \item Frames with conflicts or with arrows only are dropped.
\end{enumerate}

\chapter{Hyperparameters}
\label{app:hparams}

\begin{table}[h]
\centering
\caption{Training settings of the direction-aware model.}
\label{tab:hparams-dam}
\small
\begin{tabularx}{\textwidth}{@{}lX@{}}
\toprule
Setting & Value \\
\midrule
Backbone & Swin-T \code{swin_tiny_patch4_window7_224} (ImageNet-1k), LayerNorm per output stage, drop-path 0.1 \\
Neck & FPN, 256 channels, P2--P5 \\
Detector & DETR, 6 encoder layers (P5), 6 decoder layers (P3, P4, encoded P5), 50 queries, 3 classes \\
ROI tokens & $K=16$, RoIAlign on P2, $7\times7$ \\
Direction attention & 4 queries, 8 heads, 1 layer + LayerNorm + FFN \\
Loss & $\lambda_c=1$, $\lambda_p=1$, $\lambda_d=0.2$; DETR: class + $5\,L_1$ + $2\,$GIoU, auxiliary losses \\
Optimizer & AdamW, lr $10^{-4}$ (backbone $10^{-5}$), weight decay $10^{-4}$ \\
Schedule & cosine over 24 epochs, 500 warm-up iterations, stopped after epoch 14 \\
Gradient clipping & 0.1 \\
Batch & 4 (one canvas shape per batch) \\
Input & long side 1024, aspect kept, padded to a multiple of 32 \\
Augmentation & brightness and contrast jitter $[0.75, 1.25]$; no flips, no hue shift \\
Precision & fp16 backbone/FPN, bf16 detector, fp32 ROI/direction branch, TF32 \\
\bottomrule
\end{tabularx}
\end{table}

\begin{table}[h]
\centering
\caption{Training settings of the structured RF-DETR baseline.}
\label{tab:hparams-rfdetr}
\small
\begin{tabularx}{\textwidth}{@{}lX@{}}
\toprule
Setting & Value \\
\midrule
Model & RF-DETR Medium, 34.6M parameters, fine-tuned from the released weights \\
Classes & \code{<moves>|<color>}: 46 (OpenLane-V2 A+B), 24 (DTLD) \\
Resolution & 1024$\times$1024 square resize, no multi-scale, no scale jitter \\
Augmentation & brightness and contrast jitter; no flips, no crops \\
Optimizer & library default; lr $10^{-4}$, encoder lr $1.5\cdot10^{-4}$, EMA weights evaluated \\
Batch & 4, gradient accumulation 4 \\
Epochs & 4 (OpenLane-V2), 15 with per-epoch report (DTLD) \\
Inference & score threshold 0.20 (OpenLane-V2, from tuning), 0.25 (DTLD, from tune split); class-agnostic NMS at IoU 0.5 \\
\bottomrule
\end{tabularx}
\end{table}

\chapter{Per-Epoch Results}
\label{app:epochs}

\begin{table}[h]
\centering
\caption{TSS per epoch on subset A, subset B and A+B val, direction-aware model with both backbones, and
subset A False-Go.}
\label{tab:epochs}
\small
\begin{tabular}{@{}r cccc cccc@{}}
\toprule
 & \multicolumn{4}{c}{Swin-T} & \multicolumn{4}{c}{ConvNeXtV2-T} \\
\cmidrule(lr){2-5}\cmidrule(lr){6-9}
Epoch & A & B & A+B & $\FG_A$ & A & B & A+B & $\FG_A$ \\
\midrule
0 & 0.891 & 0.910 & 0.899 & 0.124 & 0.863 & 0.886 & 0.872 & 0.169 \\
1 & 0.893 & 0.895 & 0.894 & 0.182 & 0.875 & 0.904 & 0.887 & 0.181 \\
2 & 0.888 & 0.925 & 0.903 & 0.223 & 0.890 & 0.916 & 0.900 & 0.150 \\
3 & 0.896 & 0.927 & 0.909 & 0.145 & 0.886 & 0.908 & 0.895 & 0.228 \\
4 & 0.904 & 0.928 & 0.914 & 0.107 & 0.884 & 0.908 & 0.894 & 0.183 \\
5 & 0.899 & 0.936 & 0.914 & 0.170 & 0.895 & 0.922 & 0.906 & 0.192 \\
6 & 0.899 & 0.926 & 0.910 & 0.139 & 0.890 & 0.905 & 0.896 & 0.222 \\
7 & 0.904 & 0.937 & 0.918 & 0.112 & 0.893 & 0.900 & 0.896 & 0.204 \\
8 & 0.899 & 0.928 & 0.911 & 0.105 & 0.896 & 0.928 & 0.909 & 0.174 \\
9 & 0.893 & 0.912 & 0.901 & 0.189 & 0.884 & 0.919 & 0.898 & 0.241 \\
10 & 0.908 & 0.927 & 0.916 & 0.132 & 0.892 & 0.914 & 0.901 & 0.186 \\
11 & 0.887 & 0.891 & 0.889 & 0.212 & 0.879 & 0.912 & 0.893 & 0.209 \\
12 & 0.903 & 0.920 & 0.910 & 0.135 & 0.891 & 0.910 & 0.899 & 0.188 \\
13 & 0.897 & 0.912 & 0.903 & 0.141 & 0.905 & 0.932 & 0.916 & 0.112 \\
14 & 0.889 & 0.911 & 0.898 & 0.164 & 0.894 & 0.919 & 0.905 & 0.194 \\
\bottomrule
\end{tabular}
\end{table}

\begin{table}[h]
\centering
\caption{RF-DETR-M on OpenLane-V2 val per epoch (threshold 0.20). The reported model is epoch 3.}
\label{tab:rfdetr-epochs}
\small
\begin{tabular}{@{}rccccc@{}}
\toprule
Epoch & TSS A & TSS B & TSS A+B & $\FG$ A+B & $\FS$ A+B \\
\midrule
0 & 0.909 & 0.919 & 0.913 & 0.113 & 0.029 \\
1 & 0.900 & 0.919 & 0.908 & 0.133 & 0.025 \\
2 & 0.902 & 0.913 & 0.907 & 0.135 & 0.021 \\
3 & 0.895 & 0.914 & 0.903 & 0.130 & 0.024 \\
\bottomrule
\end{tabular}
\end{table}

\chapter{Reproducing the Results}
\label{app:repro}

All commands run from the repository root. Data paths assume OpenLane-V2 under
\code{data/OpenLane-V2/} and DTLD, LISA and ATLAS under the directory given by the \code{DATASETS}
environment variable.

\begin{lstlisting}
# OpenLane-V2 labels (per subset and split)
python -m dtsu.labels --root data/OpenLane-V2/subset_A --split train \
  --data-dict data/OpenLane-V2/subset_A/data_dict_subset_A.json \
  --out data/labels/subset_A_train.jsonl
# DTLD and LISA labels, fit/tune splits
bash tools/build_ext_labels.sh
# unit tests (labels, metrics)
python -m pytest tests
# direction-aware model, Swin-T (resumable) and per-direction evaluation
bash tools/run_full.sh
python tools/eval_directions.py
# backbone ablation
bash tools/run_convnextv2.sh
# RF-DETR: tuning, final OpenLane-V2 run, DTLD run
bash tools/run_rfdetr_tuning.sh
bash tools/run_rfdetr.sh
bash tools/run_rfdetr_dtld.sh
# paired segment-level bootstrap
python tools/bootstrap_compare.py
# out-of-distribution evaluation of the DTLD model
python tools/eval_rfdetr_ood.py --target lisa --epoch 4 --source-threshold 0.25
python tools/eval_rfdetr_ood.py --target openlane --epoch 4 --source-threshold 0.25
python tools/eval_rfdetr_atlas.py --epoch 4 --threshold 0.25
# figures of this thesis
python thesis/make_figures.py
\end{lstlisting}

The main outputs are:
\begin{itemize}
  \item \code{runs/swin_t_AB/}: checkpoint, final and per-epoch metrics, per-direction results;
  \item \code{runs/convnextv2_t_AB/}: backbone ablation;
  \item \code{runs/rfdetr_m_AB/}, \code{runs/rfdetr_m_dtld/}: structured baseline, including
        \code{ood_lisa/}, \code{ood_openlane/} and \code{ood_atlas/};
  \item \code{runs/bootstrap_swin_vs_rfdetr.json}: bootstrap comparison.
\end{itemize}
